% ------------------------------------------------------------------
%  Capitolo 7 — Perché ci inganniamo
%  Parte II
%  Ricerca di riferimento: ricerca 01 §7; 03 §2.3
%  Voci bibliografiche principali: koriat2005, bjork2013, deslauriers2019, nelson1991, kornell2009, rozenblit2002
%  Epigrafe: ✔ verificata sul testo (vedi ricerca 03 e registro, ricerca 00)
%  Scheda del capitolo: ricerca/06_indice-ragionato.md, capitolo 7
% ------------------------------------------------------------------
\chapter{Perché ci inganniamo}
\label{cap:inganni}

\epigrafe[lui crede di sapere qualcosa pur non sapendolo]{οὗτος μὲν οἴεταί τι εἰδέναι οὐκ εἰδώς}{Platone, Apologia di Socrate 21d}

\iniziale{Q}{uesto capitolo} \segnaposto{apertura: da scrivere — vedi la scheda nell'indice ragionato}.

\section{La sensazione di sapere}

\segnaposto{da scrivere}

\section{L'esperimento di Harvard}

\segnaposto{da scrivere}

\section{Le illusioni dello studente}

\segnaposto{da scrivere}

\section{Come giudicare onestamente ciò che si sa}

\segnaposto{da scrivere}

\begin{daricordare}
\segnaposto{il principio del capitolo in due righe}
\end{daricordare}

\begin{domande}
  \item \segnaposto{domanda di recupero su questo capitolo}
  \item \segnaposto{domanda di applicazione a una materia che studi}
  \item \segnaposto{domanda di ripresa su un capitolo precedente}
\end{domande}

\begin{sintesi}
  \item \segnaposto{punto di sintesi}
\end{sintesi}

\finecapitolo
